\begin{proof}
At $T=\infty$, $d_\infty=d_{CE}+Q$, and $d_{CE}$ vanishes on arity zero; by (ii) and (iii) the arity-one component along $\Delta$ of $d_\infty O[\infty]=L(O)[\infty]$ is $Q$ of the arity-one part of $O[\infty]$ along $\Delta$, which is $\big(\{S^{(1)}(\Delta),O\}+D_1(O)\big)[\infty]$; that arity-one part is the $\varepsilon\to0$ limit of the $\delta$-part of the linearised flow, at $J=I-I^{CT}[\varepsilon]-\hbar F_1^{(0)}+\delta\,O(\varepsilon)$, of the arity-one vertex of Definition \ref{def:c9-data} minus $\delta$ times the arity-one counterterm of $O$. The graph sums converge at $T=\infty$ at generic momenta, a composite line being of order $|p|^{-3}$ at small momentum and the one-loop integrals with external legs at generic momenta converging (Propositions \ref{prop:c8-Tsafe}, \ref{prop:c9-composite} (iv)); by Theorem \ref{thm:c8-L} (ii) the local representative is the same at every finite $T$.

(1) The vertex $S^{(1)}(\Delta)$. By \eqref{eq:c9-D1b} at $T=\infty$, $Q$ of its linearised flow is $\{S^{(1)}(\Delta),W\}_\varepsilon$ minus the terms with $QW$. The special-edge part of $\{S^{(1)}(\Delta),W\}_\varepsilon$ on external legs carries $e^{-\varepsilon|p|^2}-1\to0$. By \eqref{eq:c9-ward}, $\{S^{(1)}(\Delta),W\}$ is the linearised flow of $(v.-4\varepsilon\frac d{d\varepsilon})J$ plus $4\varepsilon\frac d{d\varepsilon}W$, which tends to $0$; $v.I=v.\hbar F_1^{(0)}=0$, minus $(v.-4\varepsilon\frac d{d\varepsilon})I^{CT}[\varepsilon]$ is $4$ times the coefficient of $\log\varepsilon$ in $I^{CT}[\varepsilon]$, and by \eqref{eq:c10-scaling} $(v.-4\varepsilon\frac d{d\varepsilon})O(\varepsilon)$ is $(v.O)(\varepsilon)$ plus $4$ times the coefficient of $\log\varepsilon$ in $\partial_\delta I^{CT}[\varepsilon](O)$. The $\delta$-part tends to $(v.O)[\infty]$ plus $4$ times that coefficient at scale $\infty$ plus $4$ times the tree insertion of the coefficient of $\log\varepsilon$ in $I^{CT}[\varepsilon]$ into $O[\infty]$ (Corollary \ref{cor:c9-CS}). In the terms with $QW$, $QW$ is the flow of the arity-zero defect; by (i) its limit vanishes and by (ii) its $\delta$-part too, and $Q$ of the composite graphs remains, since $-\tfrac\hbar2\partial(QW_{\rm tree})=Q(\tfrac\hbar2\partial W_{\rm tree})$ for the operator $\partial$ of the composite kernel (Proposition \ref{prop:c9-composite} (iii)).

(2) The vertex $-\hbar F_1^{(1)}(\Delta)$ enters by tree insertion, and by Lemma \ref{lem:c6-bvexact} with $(Q+\{I,-\})O=0$, $-Q$ of its tree insertion into $O$ is minus the tree insertion of $(Q+\{I,-\})\hbar F_1^{(1)}(\Delta)$, which is $4$ times the coefficient of $\log\varepsilon$ in $I^{CT}[\varepsilon]$ plus the composite part, minus $\{\hbar F_1^{(1)}(\Delta),O\}[\infty]$; the first term cancels the last term of (1).

(3) The counterterm vertices remove the singular parts of the composite graphs (Proposition \ref{prop:c9-structure} (i), Definition \ref{def:c8-quant}); their images under $Q$ complete $Q$ of the composite graphs to $Q$ of the renormalised ones.

Collecting, $\big(\{S^{(1)}(\Delta),O\}+D_1(O)\big)[\infty]$ is $(v.O)[\infty]$ plus $4$ times the coefficient of $\log\varepsilon$ in $\partial_\delta I^{CT}[\varepsilon](O)$, minus $\{\hbar F_1^{(1)}(\Delta),O\}$, plus the composite term, all at scale $\infty$, with $v.O=\{S^{(1)}(\Delta),O\}$; by (iii) and Theorem \ref{thm:c8-L} (i) this is \eqref{eq:c10-D1}, and the local representative is the one given by the master formula, both being the unique local functional with this tree-level flow.

Letters. Four times the coefficient of $\log\varepsilon$ is two-body by Theorem \ref{thm:c10-decomp}; the brackets with $\hbar F_1^{(1)}(\Delta)$ and with the field-counting functional are derivations; composite graphs of the first kind use two letters, of the second kind one. For the second kind the identity of Lemma \ref{lem:c6-bvexact} is used for tree insertions at generic momenta, where its proof, the derivation property of $Q$ and $QP\star+P\star Q=1$ on the tree edges at nonzero momentum, applies to the nonlocal composite two-point functional and gives its bracket with $O$, nonlocal unless the functional is local in the sector of the letter. On $X=\partial^\alpha A|_{W_{\mathbb R}^\perp}$: $\{\int(A,A^*),X\}=X$, the bracket of the sector part with $X$ is $v.X$, and $-8X$ from the weight $2$, $+\tfrac23X$ from the translation-invariant solution, $-\theta\,v.X$ from \eqref{eq:c9-Fperp} and $\tfrac53v.X-\tfrac23X$ from Remark \ref{rem:c9-Pi} sum to $-8X+(\tfrac53-\theta)v.X$, in units. The table: the weights $2$, $4$, $1$ of $A$, $\lambda$, $c$, and $\{\int(A,A^*),X\}=X$, $\{\int(\lambda^*,\lambda),\lambda\}=-\lambda$, $\{\int(c^*,c),c\}=-c$ (Lemma \ref{lem:c4-count}): $-8+\tfrac23=-\tfrac{22}3$, $-16+\tfrac13=-\tfrac{47}3$, $-4+\tfrac43=-\tfrac83$.
\end{proof}
